%% =====================================================================
%%  3 / 7   Methodology  方法                                    成员 3
%%  仿写对象：课本 Unit 4 p.57--58 范文（三段，每段一类信息，细节放括号）、
%%            p.60 被动语态例句、p.68--70 方法描述要点。
%%  段落分工：Overview；Phase 1 的对象/抽样 + 步骤；Phase 2 的对象/材料 + 步骤/统计。
%%  写法：被动 + 过去时（已做）；将做的步骤用 will be；Firstly/Secondly/After；
%%        数字、单位、软件、缩写放括号；不用 we，不讲“本节如何组织”。
%% =====================================================================
\section{Methodology}
\label{sec:method}

The study aims to examine under what conditions clinical action can be legitimately delegated to a medical AI agent. A qualitative review of the literature (Phase~1) was used to identify the main themes from secondary data, and a quantitative vignette experiment (Phase~2) is designed to test these themes with primary data collected by questionnaire.

\subsection{Phase 1: Purposive Review}
\label{sec:review}

The materials were seventeen sources (nine conceptual works, six benchmark evaluations, one clinical observational study and one attitude survey) published between 2019 and 2026. They were selected purposively from the literature on medical AI, so that the corpus covered all levels of autonomy, from decision support (L0) to supervised agents (L3). A source was retained if it concerned an AI system in clinical medicine, discussed the human role, oversight, accountability or trust, and was peer-reviewed or a preprint with a documented method. Sources that evaluated single-turn question answering without any action, or lay outside clinical medicine, were excluded. Foundational work on automation \citep{bainbridge1983,parasuraman2010} was used as background and was not coded.

Firstly, each source was read in the sections concerning the human role. Secondly, it was coded on four fields: the core claim, the locus of action (who acts on the patient or the record), the level of autonomy (L0--L3; \citealp{ye2025autonomy}) and the type of evidence (benchmark, clinical, attitude survey or conceptual). After coding, the core claims were grouped into themes, and a theme was retained if it was supported by at least two sources. The coded sources are listed in Table~\ref{tab:corpus}.

\subsection{Phase 2: Vignette Experiment}
\label{sec:protocol}

A $2\times2\times4$ mixed design will be used. Fifty-two adult volunteers (aged 18--45 years) without a clinical qualification will be recruited through university networks and online groups, and randomly assigned to four groups of 13 according to the role of the system (advisor or agent) and the disclosure of its identity (AI agent or unspecified). A convenience sample will be used because of limited time and resources. Four fictional vignettes (80--120 words each) will be written to span clinical risk: an oral hypoglycaemic adjustment (low), the withholding of an ACE inhibitor for hyperkalaemia (moderate), the automatic escalation of chest-pain triage (moderate to high) and an automated adverse-drug-reaction notice sent to a patient (high). Each vignette will be prepared in four versions, which differ only in the final sentence.

Firstly, the four vignettes will be presented to each participant in a random order. Secondly, after each vignette, trust (three items), willingness to delegate (two items, one reverse-scored), responsibility attribution (one forced-choice item) and perceived risk (one item) will be measured on seven-point scales, and a manipulation check will ask who completed the action. Finally, background information (age, gender and field of study) will be collected. Responses completed in less than 60 seconds or failing the manipulation check will be excluded. The trust items will be averaged once Cronbach's $\alpha$ exceeds 0.7, and the ratings will be analysed with a mixed-design analysis of variance (ANOVA) using JASP. Responsibility attribution will be tested with a chi-square test, and $p\leq0.05$ will be considered statistically significant. If fewer than 40 valid responses are obtained, the results will be reported descriptively.
